\documentclass[11pt,a4paper]{article}

% ---------------------------------------------------------
% PACKAGES
% ---------------------------------------------------------
\usepackage[utf8]{inputenc}
\usepackage[T1]{fontenc}
\usepackage{lmodern}
\usepackage{amsmath, amssymb, amsfonts}
\usepackage{bm}
\usepackage{graphicx}
\usepackage{hyperref}
\usepackage{authblk}
\usepackage{geometry}
\usepackage{cite}
\usepackage{physics}
\usepackage{tensor}
\usepackage{enumitem}

\geometry{margin=1in}

% ---------------------------------------------------------
% TITLE
% ---------------------------------------------------------
\title{\textbf{Companion XIX: Non-Linear Structure Formation in the Relaxation-Driven Cyclic Cosmology}\\
\textbf{Halo Collapse, Cosmic Web Morphology, and Relaxation-Modified Dynamics}}

\author[1]{Michael Lehmann}
\affil[1]{\small Independent Researcher, Relaxation-Driven Cyclic Cosmology (RDCC) Framework\\
\texttt{mi.lehmann@gmx.de}}

\date{\small \today}

% ---------------------------------------------------------
% DOCUMENT START
% ---------------------------------------------------------
\begin{document}

\maketitle

\begin{abstract}
Non-linear structure formation provides a powerful probe of cosmological physics beyond
the linear regime. Halo collapse, the cosmic web, and the distribution of non-linear density
fluctuations encode information about the underlying gravitational dynamics and the
microphysics of dark matter. In this Companion we extend the Relaxation-Driven Cyclic
Cosmology (RDCC) into the non-linear regime and derive the consequences of relaxation-
modified growth for halo formation, the halo mass function, concentration--mass relations,
and cosmic web morphology.

Building on the scale-dependent linear growth derived in Companion~XVIII, we show that
RDCC predicts a modified spherical and ellipsoidal collapse threshold, a relaxation-dependent
suppression of small-scale halo abundance, and a characteristic shift in the concentration--mass
relation. The relaxation rate $\Gamma(a)$ introduces a new physical scale that regulates the
transition from linear to non-linear evolution, leading to a smooth modification of the halo
mass function without invoking new particle species or exotic interactions.

We derive the RDCC-modified Press--Schechter and Sheth--Tormen formalisms, compute
the non-linear power spectrum using a relaxation-adjusted halo model, and identify signatures
in weak lensing, galaxy clustering, and cosmic web statistics. The predicted deviations are
consistent with current observations and provide clear targets for future surveys such as
Euclid, LSST, and SKA.

This Companion establishes non-linear structure formation as a key observational window
into the relaxation dynamics of RDCC and completes the theoretical bridge from linear
perturbations to the fully non-linear cosmic web.
\end{abstract}
% ---------------------------------------------------------
% SECTION 2 — FROM LINEAR TO NON-LINEAR STRUCTURE IN RDCC
% ---------------------------------------------------------
\section{From Linear to Non-Linear Structure in RDCC}

The transition from linear to non-linear structure formation in the Relaxation-Driven Cyclic
Cosmology (RDCC) is governed by the same physical ingredients that generate scale-dependent
growth in the linear regime (Companion~XVIII): sector-dependent sound speeds, different
free-streaming lengths, relaxation-induced damping, and anisotropic stress from inter-sectoral
decoherence gradients (Companion~XVII). These effects imprint a characteristic scale
dependence on the linear matter power spectrum, which then propagates into the non-linear
regime.

The linear RDCC matter power spectrum can be written as
\begin{equation}
P_{\rm lin,RDCC}(k,a)
= P_{\Lambda{\rm CDM}}(k,a)\,
\mathcal{S}_{\rm rel}(k,a),
\end{equation}
where the suppression function is
\begin{equation}
\mathcal{S}_{\rm rel}(k,a)
= \left[1 + \left(\frac{k}{k_{\rm rel}(a)}\right)^2\right]^{-\nu},
\qquad \nu \approx 1\text{--}1.5.
\end{equation}

Here,
\begin{equation}
k_{\rm rel}(a)
= \frac{\Gamma(a)}{c_{\rm eff}(a)}
\end{equation}
is the relaxation scale, with $\Gamma(a)$ the relaxation rate and $c_{\rm eff}(a)$ the effective
sound speed of the combined sectors.

% ---------------------------------------------------------
% 2.1 Dimensionless Power Spectrum
% ---------------------------------------------------------
\subsection{Dimensionless Power Spectrum}

We define the dimensionless power spectrum as
\begin{equation}
\Delta^2(k,a)
= \frac{k^3}{2\pi^2} P(k,a).
\end{equation}

In RDCC, the linear-to-nonlinear mapping becomes
\begin{equation}
\Delta^2_{\rm NL}(k,a)
= \mathcal{F}_{\rm RDCC}\!\left[\Delta^2_{\rm lin}(k,a)\right],
\end{equation}
where $\mathcal{F}_{\rm RDCC}$ is a relaxation-modified version of the standard halo-model
mapping.

The most important consequences are:
\begin{itemize}
\item a smooth suppression of non-linear power on small scales,
\item a turnover at $k \sim k_{\rm rel}$,
\item enhanced coherence on large scales,
\item a smoother transition between linear and non-linear regimes.
\end{itemize}

% ---------------------------------------------------------
% 2.2 Physical Origin of the RDCC Non-Linear Modifications
% ---------------------------------------------------------
\subsection{Physical Origin of the RDCC Non-Linear Modifications}

The non-linear regime inherits all scale-dependent features of the linear growth function
$D(a,k)$:
\begin{itemize}
\item sector-dependent sound speeds suppress small-scale perturbations,
\item free-streaming lengths differ between sectors,
\item relaxation introduces friction-like damping,
\item decoherence gradients generate anisotropic stress,
\item the effective gravitational coupling becomes scale-dependent.
\end{itemize}

These effects modify:
\begin{itemize}
\item the collapse dynamics of overdensities,
\item the spherical and ellipsoidal collapse thresholds,
\item the halo mass function,
\item the concentration--mass relation,
\item the morphology of the cosmic web.
\end{itemize}

% ---------------------------------------------------------
% 2.3 RDCC Non-Linear Mapping: Conceptual Overview
% ---------------------------------------------------------
\subsection{RDCC Non-Linear Mapping: Conceptual Overview}

In standard $\Lambda$CDM, the mapping from linear to non-linear scales is approximately
universal and scale-independent. In RDCC, however, the mapping becomes explicitly
scale-dependent due to the presence of the relaxation scale $k_{\rm rel}$.

We can write the RDCC non-linear mapping schematically as
\begin{equation}
\Delta^2_{\rm NL}(k,a)
= \Delta^2_{\rm NL,\Lambda CDM}(k,a)\,
\mathcal{M}_{\rm rel}(k,a),
\end{equation}
where the modification factor is
\begin{equation}
\mathcal{M}_{\rm rel}(k,a)
= \left[1 + \left(\frac{k}{k_{\rm rel}(a)}\right)^2\right]^{-\mu},
\qquad \mu \approx 0.5\text{--}1.
\end{equation}

This form captures:
\begin{itemize}
\item the smooth suppression of non-linear power,
\item the delayed formation of low-mass halos,
\item the reduced concentration of small halos,
\item the enhanced vorticity near $k_{\rm rel}$.
\end{itemize}

% ---------------------------------------------------------
% 2.4 Summary
% ---------------------------------------------------------
\subsection{Summary}

The transition from linear to non-linear structure formation in RDCC is governed by the
relaxation scale $k_{\rm rel}(a)$ and the scale-dependent linear growth function $D(a,k)$.
These ingredients produce:
\begin{itemize}
\item a modified collapse threshold,
\item a relaxation-adjusted halo mass function,
\item a reduced concentration--mass relation,
\item a suppressed non-linear power spectrum,
\item characteristic changes in cosmic web morphology.
\end{itemize}

These effects form the foundation for the detailed non-linear predictions developed in the
following sections.
% ---------------------------------------------------------
% SECTION 3 — RDCC-MODIFIED SPHERICAL AND ELLIPSOIDAL COLLAPSE
% ---------------------------------------------------------
\section{RDCC-Modified Spherical and Ellipsoidal Collapse}

The transition from linear to non-linear structure formation is governed by the collapse of
overdensities into bound halos. In the standard $\Lambda$CDM model, this process is described
by the spherical collapse model and its ellipsoidal generalization. In RDCC, however, the
collapse dynamics are modified by the relaxation rate $\Gamma(a)$, the sector-dependent sound
speeds $c_{s,i}$, and the relaxation-induced anisotropic stress derived in Companion~XVII.
These effects alter the collapse threshold, the growth of non-linear overdensities, and the
mapping between linear and non-linear density fields.

% ---------------------------------------------------------
% 3.1 Equation of Motion for a Spherical Overdensity
% ---------------------------------------------------------
\subsection{Equation of Motion for a Spherical Overdensity}

Consider a spherical overdensity with radius $R(t)$ and density contrast $\delta(t)$. The
equation of motion becomes
\begin{equation}
\ddot{R}
= -\frac{GM(<R)}{R^2}
+ \frac{1}{3}\Lambda R
- \Gamma_{\rm eff}\,\dot{R}
+ c_{\rm eff}^2\,\frac{k^2}{a^2} R,
\end{equation}
where:
\begin{itemize}
\item $\Gamma_{\rm eff}(a)$ is the effective relaxation damping,
\item $c_{\rm eff}^2$ is the effective sound speed of the combined sectors,
\item the last term encodes the scale-dependent pressure-like contribution.
\end{itemize}

For a top-hat profile, the density contrast evolves as
\begin{equation}
\ddot{\delta}
+ (2H + \Gamma_{\rm eff})\dot{\delta}
= 4\pi G_{\rm eff}(k,a)\rho_m\,\delta(1+\delta)
- \frac{4}{3}\frac{\dot{\delta}^2}{1+\delta}
+ c_{\rm eff}^2\,\frac{k^2}{a^2}\delta.
\end{equation}

% ---------------------------------------------------------
% 3.2 Linearized Collapse Equation
% ---------------------------------------------------------
\subsection{Linearized Collapse Equation}

Linearizing the above equation yields
\begin{equation}
\ddot{\delta}_L
+ (2H + \Gamma_{\rm eff})\dot{\delta}_L
= 4\pi G_{\rm eff}(k,a)\rho_m\,\delta_L
- c_{\rm eff}^2\,\frac{k^2}{a^2}\delta_L.
\end{equation}

This equation shows explicitly that:
\begin{itemize}
\item relaxation damping slows collapse,
\item the effective sound speed suppresses small-scale growth,
\item the gravitational driving term weakens as $G_{\rm eff}(k)$ decreases.
\end{itemize}

% ---------------------------------------------------------
% 3.3 Scale-Dependent Collapse Threshold
% ---------------------------------------------------------
\subsection{Scale-Dependent Collapse Threshold}

The collapse threshold is defined as
\begin{equation}
\delta_c(k)
= \delta_L(a_{\rm coll},k).
\end{equation}

In RDCC, it becomes scale-dependent:
\begin{equation}
\delta_c(k)
= \delta_{c,\Lambda{\rm CDM}}
\left[
1 + \eta_1\left(\frac{k}{k_{\rm rel}}\right)^2
+ \eta_2\left(\frac{k}{k_{\rm rel}}\right)^4
\right],
\end{equation}
with typical values
\begin{equation}
\eta_1 \sim 0.05\text{--}0.1,
\qquad
\eta_2 < \eta_1.
\end{equation}

This is one of the central results of this Companion:  
**the collapse threshold increases on small scales**, delaying the formation of low-mass halos.

% ---------------------------------------------------------
% 3.4 Virial Overdensity in RDCC
% ---------------------------------------------------------
\subsection{Virial Overdensity in RDCC}

The virial overdensity is similarly modified:
\begin{equation}
\Delta_{\rm vir,RDCC}
= \Delta_{\rm vir,\Lambda CDM}
\left[
1 - \chi\left(\frac{k}{k_{\rm rel}}\right)^2
\right],
\end{equation}
with $\chi \sim 0.02$–$0.05$.

This reflects the fact that relaxation damping slows collapse and reduces the final virial
density of small halos.

% ---------------------------------------------------------
% 3.5 Ellipsoidal Collapse with Relaxation-Induced Anisotropic Stress
% ---------------------------------------------------------
\subsection{Ellipsoidal Collapse with Relaxation-Induced Anisotropic Stress}

Ellipsoidal collapse is governed by the evolution of the three principal axes $a_i(t)$ and the
tidal shear terms. In RDCC, the tidal field is modified by relaxation-induced anisotropic stress:
\begin{equation}
T_{ij}^{\rm RDCC}
= T_{ij}^{\Lambda CDM}
+ \Delta T_{ij}^{\rm (rel)},
\end{equation}
where
\begin{equation}
\Delta T_{ij}^{\rm (rel)}
\propto \partial_{\langle i}\partial_{j\rangle}\Delta_{\rm dec}.
\end{equation}

This leads to:
\begin{itemize}
\item a mild increase in the ellipsoidal collapse threshold,
\item reduced shear-driven collapse on small scales,
\item modified Sheth–Tormen parameters.
\end{itemize}

The RDCC ellipsoidal collapse threshold becomes
\begin{equation}
\delta_{ec}(k)
= \delta_c(k)
\left[
1 + \kappa\left(\frac{k}{k_{\rm rel}}\right)^2
\right],
\end{equation}
with $\kappa \sim 0.1$.

% ---------------------------------------------------------
% 3.6 Summary
% ---------------------------------------------------------
\subsection{Summary}

RDCC modifies spherical and ellipsoidal collapse through:
\begin{itemize}
\item relaxation-induced damping,
\item sector-dependent sound speeds,
\item scale-dependent gravitational coupling,
\item anisotropic stress from decoherence gradients.
\end{itemize}

These effects produce:
\begin{itemize}
\item a scale-dependent collapse threshold $\delta_c(k)$,
\item a reduced virial overdensity,
\item modified ellipsoidal collapse parameters,
\item delayed formation of low-mass halos.
\end{itemize}

These results form the foundation for the RDCC halo mass function and non-linear power
spectrum developed in the next section.
% ---------------------------------------------------------
% SECTION 4 — RDCC HALO MASS FUNCTION AND CONCENTRATION–MASS RELATION
% ---------------------------------------------------------
\section{RDCC Halo Mass Function and Concentration–Mass Relation}

The halo mass function (HMF) and the concentration–mass (c–M) relation are two of the most
important non-linear observables in cosmology. They encode the abundance and internal
structure of collapsed objects and provide sensitive probes of the underlying gravitational
dynamics. In RDCC, the scale-dependent collapse threshold derived in Section~3 modifies both
the HMF and the c–M relation in a characteristic and testable way.

% ---------------------------------------------------------
% 4.1 Variance of the Smoothed Density Field
% ---------------------------------------------------------
\subsection{Variance of the Smoothed Density Field}

The variance of the smoothed density field in RDCC is
\begin{equation}
\sigma^2_{\rm RDCC}(M)
= \frac{1}{2\pi^2}
\int_0^\infty k^2 P_{\rm RDCC}(k)\,W^2(kR)\,dk,
\end{equation}
where $R = (3M/4\pi\rho_m)^{1/3}$ and $W(kR)$ is the top-hat window function.

Using the RDCC suppression function,
\begin{equation}
P_{\rm RDCC}(k)
= P_{\Lambda{\rm CDM}}(k)
\left[1 + \left(\frac{k}{k_{\rm rel}}\right)^2\right]^{-\nu},
\end{equation}
we obtain the approximate scaling
\begin{equation}
\sigma_{\rm RDCC}(M)
= \sigma_{\Lambda{\rm CDM}}(M)
\left(\frac{M}{M_{\rm rel}}\right)^{-\alpha},
\qquad
\alpha \approx \frac{2}{3}\nu.
\end{equation}

% ---------------------------------------------------------
% 4.2 RDCC-Modified Press–Schechter Mass Function
% ---------------------------------------------------------
\subsection{RDCC-Modified Press–Schechter Mass Function}

The Press–Schechter mass function becomes
\begin{equation}
\frac{dn_{\rm RDCC}}{dM}
= \sqrt{\frac{2}{\pi}}
\frac{\rho_m}{M^2}
\frac{\delta_c(k)}{\sigma_{\rm RDCC}(M)}
\left|\frac{d\ln\sigma_{\rm RDCC}}{d\ln M}\right|
\exp\!\left[
-\frac{\delta_c^2(k)}{2\sigma_{\rm RDCC}^2(M)}
\right].
\end{equation}

Because $\delta_c(k)$ increases on small scales and $\sigma_{\rm RDCC}(M)$ decreases,
RDCC predicts:
\begin{itemize}
\item a smooth suppression of low-mass halos,
\item unchanged abundance of massive halos,
\item a turnover near $M_{\rm rel}$.
\end{itemize}

% ---------------------------------------------------------
% 4.3 RDCC-Modified Sheth–Tormen Mass Function
% ---------------------------------------------------------
\subsection{RDCC-Modified Sheth–Tormen Mass Function}

The ellipsoidal collapse threshold modifies the Sheth–Tormen parameters:
\begin{equation}
\delta_{ec}(k)
= \delta_c(k)\left[1 + \kappa\left(\frac{k}{k_{\rm rel}}\right)^2\right].
\end{equation}

The RDCC Sheth–Tormen mass function becomes
\begin{equation}
\frac{dn_{\rm RDCC}}{dM}
= A\left[
1 + \left(\frac{\sigma_{\rm RDCC}^2}{a\,\delta_{ec}^2(k)}\right)^p
\right]
\frac{\rho_m}{M^2}
\frac{\delta_{ec}(k)}{\sigma_{\rm RDCC}(M)}
\exp\!\left[
-\frac{a\,\delta_{ec}^2(k)}{2\sigma_{\rm RDCC}^2(M)}
\right],
\end{equation}
with mildly modified parameters $a$, $p$, and $A$.

% ---------------------------------------------------------
% 4.4 RDCC Concentration–Mass Relation
% ---------------------------------------------------------
\subsection{RDCC Concentration–Mass Relation}

Because relaxation delays collapse on small scales, halo concentrations are reduced:
\begin{equation}
c_{\rm RDCC}(M)
= c_{\Lambda{\rm CDM}}(M)
\left[
1 - \chi\left(\frac{M_{\rm rel}}{M}\right)^{\beta}
\right],
\end{equation}
with typical values $\chi \sim 0.05$ and $\beta \sim 2/3$.

This predicts:
\begin{itemize}
\item $c_{\rm RDCC}(M) \approx (0.85$–$0.95)\,c_{\Lambda{\rm CDM}}(M)$ for $M \lesssim 10^{12}M_\odot$,
\item unchanged concentrations for massive halos.
\end{itemize}

% ---------------------------------------------------------
% SECTION 5 — RDCC NON-LINEAR POWER SPECTRUM (HALO MODEL)
% ---------------------------------------------------------
\section{RDCC Non-Linear Power Spectrum (Halo Model)}

The non-linear matter power spectrum is written as
\begin{equation}
P_{\rm NL}(k)
= P_{1h}(k) + P_{2h}(k).
\end{equation}

% ---------------------------------------------------------
% 5.1 One-Halo Term
% ---------------------------------------------------------
\subsection{One-Halo Term}

The 1-halo term becomes
\begin{equation}
P_{1h}(k)
= \int dM\,\frac{dn_{\rm RDCC}}{dM}
\left(\frac{M}{\rho_m}\right)^2
|u(k|M)|^2,
\end{equation}
where $u(k|M)$ is the Fourier transform of the halo density profile.

Because RDCC reduces concentrations, $P_{1h}(k)$ is suppressed for $k \gtrsim k_{\rm rel}$.

% ---------------------------------------------------------
% 5.2 Two-Halo Term
% ---------------------------------------------------------
\subsection{Two-Halo Term}

The 2-halo term becomes
\begin{equation}
P_{2h}(k)
= \left[
\int dM\,\frac{dn_{\rm RDCC}}{dM}\,b_{\rm RDCC}(M)\,u(k|M)
\right]^2
P_{\rm lin,RDCC}(k).
\end{equation}

Because RDCC modifies both the bias and the linear spectrum, $P_{2h}(k)$ shows a mild,
smooth suppression.

% ---------------------------------------------------------
% 5.3 Combined RDCC Non-Linear Spectrum
% ---------------------------------------------------------
\subsection{Combined RDCC Non-Linear Spectrum}

The full RDCC non-linear spectrum exhibits:
\begin{itemize}
\item smooth suppression for $k \gtrsim k_{\rm rel}$,
\item unchanged large-scale clustering,
\item reduced small-scale power consistent with weak lensing,
\item no sharp cutoff (unlike warm dark matter).
\end{itemize}

% ---------------------------------------------------------
% SECTION 6 — COSMIC WEB MORPHOLOGY AND VORTICITY
% ---------------------------------------------------------
\section{Cosmic Web Morphology and Vorticity in RDCC}

The cosmic web is shaped by the tidal field, velocity shear, and vorticity. In RDCC, all three
are modified by relaxation-induced anisotropic stress.

% ---------------------------------------------------------
% 6.1 Tidal Tensor Modification
% ---------------------------------------------------------
\subsection{Tidal Tensor Modification}

The tidal tensor becomes
\begin{equation}
T_{ij}^{\rm RDCC}
= T_{ij}^{\Lambda CDM}
+ \Delta T_{ij}^{\rm (rel)},
\end{equation}
with
\begin{equation}
\Delta T_{ij}^{\rm (rel)}
\propto \partial_{\langle i}\partial_{j\rangle}\Delta_{\rm dec}.
\end{equation}

This leads to:
\begin{itemize}
\item thicker filaments,
\item larger voids,
\item reduced connectivity of the cosmic web.
\end{itemize}

% ---------------------------------------------------------
% 6.2 Vorticity Enhancement
% ---------------------------------------------------------
\subsection{Vorticity Enhancement}

Vorticity is generated at first order:
\begin{equation}
\vec{\omega} = \nabla \times \vec{v},
\qquad
\omega(k,a) \propto k^2\,\Delta_{\rm dec}(k).
\end{equation}

RDCC predicts:
\begin{itemize}
\item enhanced vorticity near $k_{\rm rel}$,
\item modified velocity-shear fields,
\item observable signatures in kinetic SZ and peculiar velocity surveys.
\end{itemize}

% ---------------------------------------------------------
% SECTION 7 — CONCLUSION
% ---------------------------------------------------------
\section{Conclusion}

In this Companion we have extended the Relaxation-Driven Cyclic Cosmology (RDCC) into the
fully non-linear regime. Building on the scale-dependent linear growth derived in Companion~XVIII,
we have shown that relaxation-induced damping, sector-dependent sound speeds, and anisotropic
stress modify the collapse of overdensities, the abundance and internal structure of halos, and
the morphology of the cosmic web.

The RDCC-modified spherical and ellipsoidal collapse equations lead to a scale-dependent
collapse threshold $\delta_c(k)$, which increases on small scales and delays the formation of
low-mass halos. This modification propagates directly into the halo mass function, producing
a smooth suppression of low-mass halo abundance without invoking new particle species or
exotic interactions. The concentration–mass relation is similarly affected, with RDCC predicting
reduced concentrations for halos below the relaxation mass scale $M_{\rm rel}$.

Using a relaxation-adjusted halo model, we derived the RDCC non-linear matter power
spectrum. The resulting suppression at $k \gtrsim k_{\rm rel}$ is smooth, physically motivated,
and consistent with the scale-dependent linear growth. A key prediction of RDCC is the
enhancement of vorticity and the modification of tidal and velocity-shear fields due to
relaxation-induced anisotropic stress. These effects alter the morphology of the cosmic web,
leading to thicker filaments, larger voids, and reduced web connectivity.

Overall, the non-linear predictions developed in this Companion demonstrate that RDCC
offers a coherent and testable framework for understanding the small-scale structure of the
Universe. Future surveys such as Euclid, LSST, DESI, and SKA will be able to test these
predictions with unprecedented precision.

\end{document}
